\documentclass{report}
\usepackage{amsmath,amssymb}
\setlength{\parindent}{0mm}
\setlength{\parskip}{1em}
\begin{document}
\begin{center}
\rule{6in}{1pt} \
{\large Hilari C. Tiedeman \\
{\bf Multilevel Schur Complement Preconditioning for Multi-Physics Simulations}}

Southern Methodist University \\ Department of Mathematics \\ PO Box 750156 \\ Dallas \\ Texas 75275-0156
\\
{\tt htiedeman@smu.edu}\\
Daniel R. Reynolds\end{center}

Advection-diffusion PDEs are prevalent in models of many physical
applications in science and engineering. In this talk, we focus on a
scalar-valued advection-diffusion-reaction equation of the form
\[
\partial_t u = \nabla\cdot(\beta u)+\mu\nabla\cdot(D\nabla u)+f,
\]
where $\beta$ is a vector-valued advection coefficient that may
depend on $u$, $D$ is a matrix of diffusion coefficients, and
$f(u,t)$ is a forcing term. Within different parameter regimes, there
exist optimally scalable solvers for problems of this type, however no
single solver yet applies well within all regimes of physical interest.

Domain decomposition methods display nearly optimal parallel
scalability within the advection-dominated regime, but typically do
not scale well for diffusion-dominated problems. The converse occurs
when using multigrid methods. Our goal is to find one method that is
scalable in both regimes, through using a hybrid approach combining
restrictive additive Schwarz (domain decomposition) and geometric
multigrid.

Our approach begins with a Schur complement formulation of the
linearized implicit system ($Au=f$) as with the FETI (Farhat and
Roux, 1991), BDD (Mandel, 1993), and BDDC (Dohrmann, 2003) algorithms, where the
unknowns $u$ are split into two sets: those residing in a domain
interior, $u_I$, and those residing at inter-processor boundaries,
$u_\Gamma$. With this decomposition, one may similarly decompose the
full Jacobian matrix into four associated blocks,
\[
A = \left[\begin{array}{cc}
A_{I,I} & A_{I,\Gamma} \\
A_{\Gamma,I} & A_{\Gamma,\Gamma} \\
\end{array}\right].
\]
Since the $A_{I,I}$ matrix is itself block-diagonal, with each block
corresponding to the Jacobian matrix dependencies within a single
processor, we can apply $A^{-1}_{I,I}$ using a standard sparse direct
solver. Elimination of this block results in a Schur complement system
for the interface nodes,
\[
\left(A_{\Gamma,\Gamma}-A_{\Gamma,I}A^{-1}_{I,I}A_{I,\Gamma}\right)u_\Gamma=g_\Gamma,
\]
where $g_\Gamma$ is formed during the first elimination step, and the
Schur complement matrix is given by $S =
A_{\Gamma,\Gamma}-A_{\Gamma,I}A^{-1}_{I,I}A_{I,\Gamma}$.

While traditional domain-decomposition methods do not solve this
global Schur complement system directly (and instead solve an
interface system based on only a small fraction of unknowns), we solve
the full interface system using a multilevel technique similar to
multigrid. Here, our fine grid problem consists of the entire
Schur complement (interface) system. We then proceed through a
traditional set of V-cycle iterations, where at each level we obtain
an increasingly coarse subset of the full interface system. However,
due to the $A^{-1}_{I,I}$ term in $S$, all residual corrections are
evaluated in a matrix-free fashion, using FGMRES as a smoother.

After presenting the details of our algorithm, we present simulation
results using the Ranger supercomputer at TACC. We investigate
problems in both the advection- and diffusion-dominated regimes, and
examine scalability of both iteration count and wall-clock time.


\end{document}
